\chapter{书系路线图}
\label{app:series-roadmap}

八卷并不是把朝代依次排开，而是追问一个问题怎样在新的尺度和材料条件下重现。第一卷从西周的封国网络写到秦的统一行政；它留下的不是一条“必然走向帝国”的直线，而是一组后继政权必须重新处理的难题：如何让中央命令抵达地方，如何取得粮食、兵员与信息，如何安排地方精英和边疆，如何使动员不至于耗尽信任。

\section{八卷的接力}

\begin{tabularx}{\linewidth}{@{}p{0.14\linewidth}p{0.22\linewidth}X@{}}
\toprule
\textbf{卷} & \textbf{时间范围} & \textbf{承接的问题} \\
\midrule
第一卷：周秦 & 西周至前206年 & 封国、会盟、县邑和帝国行政如何先后解决距离、动员与合法性，又如何把代价留给后来者。 \\
第二卷：汉帝国 & 前206—220 & 汉如何接续并调适秦的行政遗产，并在财政、边疆、地方与精英关系中重估统一的成本。 \\
第三卷：三国 & 220—280 & 区域国家如何把战争、屯田、道路与补给组织为政权能力，并在竞争中走向再统一。 \\
第四卷：两晋南北朝 & 280—589 & 门阀、迁徙、部族政权与南北分治如何重组地方社会、军镇和统一的条件。 \\
第五卷：隋唐五代 & 581—960 & 再统一怎样依靠制度、交通与财政；藩镇、区域社会和财政转型为何又使统一难以维持。 \\
第六卷：宋辽夏金元 & 960—1368 & 多政权竞争、商业财政和跨区域网络怎样改变国家能力；征服与交流如何同时重塑边疆与社会。 \\
第七卷：明 & 1368—1644 & 皇权、赋役、白银、军防与地方社会怎样互相制约，并在晚期危机中暴露新的脆弱处。 \\
第八卷：清至1912 & 1644—1912 & 征服王朝的整合、人口与财政压力、世界性冲击与革命怎样共同终结帝国。 \\
\bottomrule
\end{tabularx}

\section{第一卷交给后卷的五条线}

\begin{description}
  \item[中央与地方] 从\eventref{WZ-EAST-ENFEOFFMENT}{西周·重组东方封国网络}到\eventref{SA-0546-MIBING}{春秋·弭兵之会}，再到\eventref{WQ-0356-SHANGYANG}{战国·商鞅变法}，中央权威总须借地方关系、盟友或官吏落地。秦的郡县与文书把这种关系改写为较直接的行政链，第二卷将考察汉代如何继承而非简单复制这条链。
  \item[动员与负担] 战国的持续竞争和\eventref{WQ-0260-CHANGPING}{战国·长平之战}显示，粮食、道路、登记和奖赏已成为战争的一部分；统一后，\eventref{QIN-0213-EMPIRE}{秦·高强度治理}又显示动员能力与社会负担不能分开。后卷将持续追踪赋役、军役、财政和市场如何相互塑造。
  \item[合法性与继承] 周王室的名分、春秋的“尊王”语言和秦的皇帝名号，都不是单纯的装饰，而是协调资源、册命和服从的方式。它们也无法在继承失序时自动生效；\eventref{QIN-0210-SUCCESSION}{秦·沙丘继承}正是下一卷理解帝国调适的起点，而不是一条隔绝的结尾。
  \item[边疆与空间] 周的东方经营、秦的关中位置、楚与吴越的水陆世界、赵燕的北方压力，提示疆界从来不是一条静止的线。各卷都将用道路、河流、城邑、迁徙和军事通道说明国家为何在不同地区采取不同的治理方式。
  \item[材料与沉默] 金文、经年、国别叙事、正史、简牍、地方文书和考古遗存将随时代改变比例。后卷不会把材料变多误作事实变得无争议；每一卷仍须区分现场文书、后出的叙事和现代解释。
\end{description}

\section{从秦末进入下一卷}

\eventref{QIN-0209-CHENSHENG}{秦·大泽起事}、\eventref{QIN-0207-JULU}{秦·巨鹿之战}与\eventref{QIN-0206-SURRENDER}{秦·子婴降汉}标记的是统一帝国的指挥链断裂、各地政治力量重组和秦政权的终止，并不自动解释汉朝随后每一项制度选择。第二卷将从这一遗产出发，分别检验哪些行政技术、精英关系、财政负担和边疆问题被接续、改造或抛弃。读者可先由\hyperref[app:index]{事件与图表索引}回看第一卷的具体场景，再用\hyperref[app:sources]{来源与史料说明}辨认这条过渡线能够说到何处。

\section{怎样跨卷而不越过材料}

第一卷的事件链接只回到本卷已经写入的锚点，不能把秦末的结果倒写成汉、三国或更晚时代已经注定的前因。读者若要追踪一条问题线，可先用\hyperref[app:index]{事件、主题、来源和图件索引}找到周秦的可见节点，再用\hyperref[app:bibliography]{基础书目与读法索引}辨认这一卷所凭的材料层次；进入后卷时，应从各卷自己的年序、来源附录和事件锚点重新开始核对。政权延续、制度沿用或同一地名，都只是值得追问的线索，不自动证明政策、社会负担或人的选择没有变化。
